% !TEX root = ../../main.tex

\section{关键用例规约}

在完成普通用户和系统管理员的功能需求分析后，为进一步明确系统关键功能的交互边界和业务约束，本文选取普通用户登录、实时手势训练、句子视频识别、手势资源管理和模型版本发布五个核心用例进行规约说明。
关键用例规约主要从参与者、前置条件、基本流程、异常流程和后置条件等方面描述系统行为，为后续系统设计与实现提供依据。

\subsection{用户登录用例规约}

用户登录是普通用户进入移动端系统并使用训练与识别功能的前提。
该用例用于验证用户身份，并为后续访问手语资源、实时训练和句子视频识别等功能提供登录状态支撑。

\begin{longtable}{p{0.20\textwidth}p{0.70\textwidth}}
  \caption{用户登录用例规约}
  \label{tab:user-login-usecase}\\
  \toprule
  \textbf{项目} & \textbf{内容} \\
  \midrule
  用例名称 & 用户登录 \\
  主要参与者 & 普通用户 \\
  前置条件 & 用户已打开移动端应用，且当前未处于有效登录状态 \\
  后置条件 & 登录成功后进入应用主界面；登录失败则停留在登录页面并给出提示 \\
  基本流程 & 1. 用户进入登录或注册页面；2. 用户输入账号信息并提交；3. 系统校验用户身份；4. 校验通过后生成登录状态；5. 系统跳转至应用主界面 \\
  异常流程 & 1. 若账号信息错误，系统提示登录失败；2. 若网络异常，系统提示用户稍后重试；3. 若登录状态过期，系统要求用户重新登录 \\
  \bottomrule
\end{longtable}

\subsection{实时手势训练用例规约}

实时手势训练是普通用户侧的核心训练功能。
用户选择手语资源后，进入对应训练页面，观看动作示范，并通过摄像头完成手势动作练习，系统向用户展示实时识别反馈。

\begin{longtable}{p{0.20\textwidth}p{0.70\textwidth}}
  \caption{实时手势训练用例规约}
  \label{tab:realtime-training-usecase}\\
  \toprule
  \textbf{项目} & \textbf{内容} \\
  \midrule
  用例名称 & 实时手势训练 \\
  主要参与者 & 普通用户 \\
  前置条件 & 用户已登录，并已选择需要练习的手语资源 \\
  后置条件 & 用户完成一次或多次手势练习，并查看实时识别结果 \\
  基本流程 & 1. 用户进入手势训练页面；2. 系统展示手语资源信息和动作示范；3. 用户开始练习；4. 用户对准摄像头完成手势动作；5. 系统展示实时识别结果；6. 用户选择继续练习或结束练习 \\
  异常流程 & 1. 若摄像头不可用，系统提示用户检查权限或设备状态；2. 若识别结果暂不可用，系统提示用户调整姿态或稍后重试；3. 若网络连接异常，系统提示识别服务连接失败 \\
  \bottomrule
\end{longtable}

\subsection{句子视频识别用例规约}

句子视频识别用于支持用户提交本地手语视频，并查看系统返回的原始识别结果和语义修正结果。
该用例体现了系统在有限词表场景下对辅助交流能力的探索。

\begin{longtable}{p{0.20\textwidth}p{0.70\textwidth}}
  \caption{句子视频识别用例规约}
  \label{tab:sentence-video-usecase}\\
  \toprule
  \textbf{项目} & \textbf{内容} \\
  \midrule
  用例名称 & 句子视频识别 \\
  主要参与者 & 普通用户 \\
  前置条件 & 用户已登录，并进入句子视频识别页面 \\
  后置条件 & 系统返回视频识别结果；用户可以查看原始识别结果和语义修正结果 \\
  基本流程 & 1. 用户进入句子视频识别页面；2. 用户选择本地手语视频；3. 用户提交识别请求；4. 系统提示正在识别；5. 识别完成后展示原始识别结果；6. 系统展示语义修正结果；7. 用户选择结束或继续识别其他视频 \\
  异常流程 & 1. 若用户未选择视频，系统提示先选择视频；2. 若视频格式或文件读取失败，系统提示文件不可用；3. 若识别失败，系统提示用户重新提交；4. 若语义修正失败，系统保留原始识别结果并提示修正暂不可用 \\
  \bottomrule
\end{longtable}

\subsection{手势资源管理用例规约}

手势资源管理是后台管理端的重要功能，用于维护移动端训练和识别所需的基础资源。
管理员通过该功能维护手势资源的基本信息、分类归属、展示资源和启用状态。

\begin{longtable}{p{0.20\textwidth}p{0.70\textwidth}}
  \caption{手势资源管理用例规约}
  \label{tab:gesture-resource-usecase}\\
  \toprule
  \textbf{项目} & \textbf{内容} \\
  \midrule
  用例名称 & 手势资源管理 \\
  主要参与者 & 系统管理员 \\
  前置条件 & 管理员已登录后台管理端，并具有资源管理权限 \\
  后置条件 & 手势资源信息被新增、修改、查询或更新状态，移动端可根据资源状态进行展示和使用 \\
  基本流程 & 1. 管理员进入手势资源管理页面；2. 管理员查询或筛选资源列表；3. 管理员新增或编辑资源信息；4. 系统校验资源字段；5. 系统保存资源信息；6. 页面刷新并展示最新资源列表 \\
  异常流程 & 1. 若必填字段缺失，系统提示补充信息；2. 若资源编码重复，系统提示编码不可重复；3. 若文件资源上传失败，系统提示重新上传；4. 若管理员未登录或权限不足，系统拒绝访问 \\
  \bottomrule
\end{longtable}

\subsection{模型版本发布用例规约}

模型版本发布用于将训练完成并登记的模型版本设置为当前可用版本，使移动端识别功能能够使用最新发布的模型能力。
该用例是后台模型维护闭环中的关键步骤。

\begin{longtable}{p{0.20\textwidth}p{0.70\textwidth}}
  \caption{模型版本发布用例规约}
  \label{tab:model-publish-usecase}\\
  \toprule
  \textbf{项目} & \textbf{内容} \\
  \midrule
  用例名称 & 模型版本发布 \\
  主要参与者 & 系统管理员 \\
  前置条件 & 管理员已登录后台管理端；系统中已存在可发布的模型版本 \\
  后置条件 & 指定模型版本被设置为当前发布版本，识别服务后续可加载并使用该模型 \\
  基本流程 & 1. 管理员进入模型版本管理页面；2. 管理员查看模型版本列表；3. 管理员选择需要发布的模型版本；4. 系统提示确认发布；5. 管理员确认发布；6. 系统更新模型发布状态；7. 系统提示发布成功 \\
  异常流程 & 1. 若模型文件缺失，系统提示该版本不可发布；2. 若发布过程失败，系统保留原发布版本并提示失败原因；3. 若管理员权限不足，系统拒绝发布操作 \\
  \bottomrule
\end{longtable}
